\begin{abstract}
Mathematical reasoning remains a formidable obstacle for language models because of its inherent complexity and structured requirements. In this work, we present DeepSeekMath~7B, an extension of DeepSeek-Coder-Base-v1.5 7B, further pre-trained with 120B mathematics-focused tokens gathered from Common Crawl, in combination with both natural language and programming code data. DeepSeekMath~7B attains a notable 51.7\% on the challenging, competition-tier MATH benchmark, reaching results close to those achieved by Gemini-Ultra and GPT-4, all without utilizing external toolkits or ensemble-based voting. Evaluating 64 generated samples via the self-consistency method, DeepSeekMath~7B obtains a 60.9\% success rate on MATH. Two principal innovations drive the mathematical reasoning skills of DeepSeekMath: (1) the exploitation of publicly available web resources by means of a carefully designed data filtration and selection strategy, and (2) the implementation of Group Relative Policy Optimization (GRPO)—a tailored version of Proximal Policy Optimization (PPO)—which augments both mathematical reasoning efficacy and the memory efficiency of PPO.
\end{abstract}

\section{Introduction}

Large language models (LLMs) have fundamentally transformed artificial intelligence approaches to mathematical reasoning, enabling considerable progress on quantitative reasoning evaluations \citep{MATH} as well as geometry reasoning tasks \citep{trinh2024solving}. These systems now serve as valuable tools in aiding users with the resolution of intricate mathematical problems \citep{tao}. Nonetheless, leading-edge solutions such as GPT-4 \citep{gpt4} and Gemini-Ultra \citep{gemini} remain inaccessible to the public, and current open-source alternatives exhibit a substantial performance gap relative to these proprietary models.

This work presents DeepSeekMath, a specialized language model designed for mathematical reasoning, which not only surpasses open-source alternatives in mathematical tasks but also reaches a proficiency comparable to GPT-4 across standard academic evaluations. Central to our approach is the development of the DeepSeekMath Corpus, a comprehensive pre-training dataset encompassing 120B math tokens. Construction of this corpus involves leveraging the Common Crawl (CC) and filtering relevant mathematical content with a classifier built upon fastText \citep{joulin2016fasttext}. The initial phase involves training the classifier on positive examples drawn from OpenWebMath \citep{openwebmath}, alongside a heterogeneous collection of non-mathematical web pages for negative sampling. Using the classifier, we then extract further mathematical content from CC, subsequently applying human annotation to ensure data quality. The improved dataset is utilized to retrain the classifier, thereby enhancing extraction precision. Our assessment shows that the resulting corpus attains high quality, as evidenced by DeepSeekMath-Base 7B achieving scores of 64.2\% on GSM8K \citep{gsm8k} and 36.2\% on the MATH competition dataset \citep{MATH}, outperforming the Minerva 540B model \citep{minerva}. Notably, since the DeepSeekMath Corpus includes multiple languages, we observe increased accuracy on Chinese mathematical evaluations \citep{wei2023cmath,agieval}. We contend that our methodology for curating mathematical datasets can inform future research, and that there exists considerable potential for subsequent advancements.

DeepSeekMath-Base is constructed by initializing with DeepSeek-Coder-Base-v1.5 7B \citep{deepseek-coder}, as empirical results suggest that leveraging a code-pretrained model yields superior results compared to models starting from general-purpose LLMs. Additionally, we find that the incorporation of mathematical pretraining not only boosts the model’s performance on mathematics-specific benchmarks but also yields gains on broader benchmarks such as MMLU \citep{mmlu} and BBH \citep{bbh}. This indicates the enhancement of both mathematical proficiency and general reasoning skills.

Subsequent to pre-training, DeepSeekMath-Base undergoes mathematical instruction fine-tuning utilizing chain-of-thought \citep{cot}, program-of-thought \citep{pot,pal}, and tool-based reasoning datasets \citep{tora}. The fine-tuned model, DeepSeekMath-Instruct 7B, surpasses all other available 7B models and reaches a competitive level with much larger, instruction-tuned models of 70B parameters.

Moreover, we present Group Relative Policy Optimization (GRPO), an adapted form of reinforcement learning (RL) rooted in Proximal Policy Optimization (PPO) \citep{schulman2017proximal}. Unlike standard PPO, GRPO omits the critic network and instead derives the baseline through group-based scoring, leading to significantly lower computational requirements. Applying GRPO on a subset of English instruction fine-tuning data, we achieve marked improvements over the already strong DeepSeekMath-Instruct, with notable gains observed in both in-domain benchmarks (GSM8K: 82.9\% $\rightarrow$ 88.2\%, MATH: 46.8\% $\rightarrow$ 51.7\%) and in out-of-domain assessments (such as CMATH: 84.6\% $\rightarrow$ 88.8\%) within the RL training phase. 

To systematize and interpret various approaches, we propose a unified paradigm encompassing methods like Rejection Sampling Fine-Tuning (RFT) \citep{yuan2023scaling}, Direct Preference Optimization (DPO) \citep{dpo}, PPO, and GRPO. Through this framework, we observe that these techniques are variations or simplifications of RL strategies. Our thorough experimental investigation—covering online versus offline RL, outcome versus process-oriented supervision, as well as single-turn against iterative RL—sheds light on the foundational aspects of this unified paradigm. Finally, we articulate the reasons behind the observed RL-induced performance gains for instruction-tuned models and outline promising avenues for more effective RL advancements within this conceptual framework.

\subsection{Contributions}
This work offers advancements in large-scale math pre-training, together with investigations into reinforcement learning approaches.

\textbf{Math Pre-Training at Scale}

\begin{itemize}[topsep=0pt]
    \item We demonstrate that Common Crawl, a freely available dataset, is a rich source of mathematical information. Through the application of a rigorous data selection pipeline, we curated the DeepSeekMath Corpus, which comprises 120B tokens extracted from mathematically relevant web pages—representing a dataset nearly 7 times larger than Minerva’s math web page collection \citep{minerva} and 9 times greater than OpenWebMath’s dataset \citep{openwebmath}.

    \item Our DeepSeekMath-Base 7B pre-trained model attains performance on par with Minerva 540B \citep{minerva}, suggesting that the quantity of parameters is not the sole determinant of mathematical reasoning performance. Rather, a compact model trained on well-curated data can achieve competitive results.

    \item We present observations from math model training, showing that performing code training before math-focused pre-training enhances a model’s capacity to tackle mathematical tasks, regardless of tool assistance. This contributes to ongoing discussions around the impact of code training on reasoning skills, particularly for math-related reasoning, and provides affirmative evidence in this context.

    \item While incorporating arXiv papers into training sets is widespread, especially for math-oriented work, we find that their inclusion does not yield discernible gains across the mathematical benchmarks considered in this study.
\end{itemize}

\textbf{Exploration and Analysis of Reinforcement Learning}

\begin{itemize}[topsep=0pt]
    \item We present Group Relative Policy Optimization (GRPO), a reinforcement learning approach that is both resource-efficient and effective. By eliminating the critic model and instead estimating baselines through group scores, GRPO markedly decreases the computational requirements in comparison to Proximal Policy Optimization (PPO).
    \item Our results show that applying GRPO notably improves the performance of the instruction-tuned model DeepSeekMath-Instruct using only instruction-tuning data. Moreover, we identify improvements in out-of-domain capabilities throughout the reinforcement learning procedure.
    \item We propose a comprehensive framework for interpreting a range of methods, including RFT, DPO, PPO, and GRPO. Additionally, we conduct a series of systematic experiments, such as online versus offline training, outcome-based versus process-based supervision, and single-turn versus iterative reinforcement learning, among others, to thoroughly analyze the fundamental components of this framework.
    \item Leveraging this unified perspective, we examine the underlying factors contributing to the success of reinforcement learning, and outline various promising avenues for enhancing reinforcement learning effectiveness in large language models (LLMs).
\end{itemize}

\subsection{Summary of Evaluations and Metrics}

\begin{itemize}[topsep=0pt]
    \item \textbf{English and Chinese Mathematical Reasoning}: We systematically evaluate our models against both English and Chinese benchmark datasets, spanning mathematical tasks from elementary to advanced collegiate levels. The English benchmarks include GSM8K \citep{gsm8k}, MATH \citep{MATH}, SAT \citep{llemma}, OCW Courses \citep{minerva}, and MMLU-STEM \citep{mmlu}; while Chinese benchmarks cover MGSM-zh \citep{mgsm}, CMATH \citep{wei2023cmath}, Gaokao-MathCloze \citep{agieval}, and Gaokao-MathQA \citep{agieval}. Evaluations encompass both the models' capacity for generating complete textual solutions unaided by external tools, and their ability to address problems with Python.

    On English benchmarks, DeepSeekMath-Base exhibits performance on par with the proprietary Minerva 540B \citep{minerva}, and consistently outperforms all open-source foundational models (such as Mistral 7B \citep{mistral} and Llemma-34B \citep{llemma}), irrespective of prior math-specific pre-training, often by a wide margin. Importantly, DeepSeekMath-Base achieves leading results on Chinese datasets, likely because our approach diverges from prior studies \citep{minerva,llemma} by supplementing English math data with substantial, high-quality non-English sources. Through instruction-based tuning and reinforcement learning, our enhanced models—DeepSeekMath-Instruct and DeepSeekMath-RL—exhibit outstanding outcomes, surpassing 50\% accuracy on the competitive MATH benchmark, which represents a first for open-source systems.
\end{itemize}

\item \textbf{Formal Mathematics}: We assess DeepSeekMath-Base’s capabilities in the informal-to-formal theorem proving task, following the framework in \citep{dsp_proof}, using the miniF2F benchmark \citep{minif2f} and selecting Isabelle \citep{isabelle} as the proof assistant. DeepSeekMath-Base shows robust performance in few-shot autoformalization scenarios.

\item \textbf{Natural Language Understanding, Reasoning, and Code}: To provide a holistic evaluation of the model’s proficiency in general understanding, reasoning, and code generation, we benchmark DeepSeekMath-Base on Massive Multitask Language Understanding (MMLU) \citep{mmlu}, which features 57 multiple-choice tasks spanning a wide range of fields; BIG-Bench Hard (BBH) \citep{bbh}, composed of 23 difficult, primarily multi-step reasoning challenges; as well as HumanEval \citep{codex} and MBPP \citep{mbpp}, both standard benchmarks for code-oriented language models. Mathematical pre-training enhances both comprehension and reasoning skills across these evaluations.

\section{Math Pre-Training}

\subsection{Data Collection and Decontamination} Here, we describe the methodology for assembling the DeepSeekMath Corpus by leveraging data from Common Crawl. As illustrated in Figure \ref{fig:data_collect}, our approach employs an iterative pipeline, systematically collecting an extensive mathematical dataset from Common Crawl, starting from a seed dataset—such as a small, high-quality set of mathematical resources. This methodology is generalizable to additional domains, including programming.

Initially, we designate OpenWebMath \citep{openwebmath}, a curated compilation of high-quality mathematical content from the web, as the starting seed corpus. With this seed, we train a fastText model \citep{joulin2016fasttext} to identify and retrieve additional mathematical webpages similar to those found in OpenWebMath. For this purpose, 500,000 items are randomly sampled from the seed corpus as positive examples, and another 500,000 webpages from Common Crawl serve as negative examples. Training leverages an open-source implementation\footnote{\url{https://fasttext.cc}}, using a vector size of 256, learning rate of 0.1, maximum word n-gram length of 3, minimum word occurrence threshold of 3, and 3 training epochs. To curtail the vastness of the original Common Crawl dataset, we apply both URL deduplication and near-duplicate filtering, reducing it to 40B HTML documents. The trained fastText model is then used to identify mathematical webpages from this deduplicated set. Subsequently, collected documents are scored using the model’s predictions, with only the highest-scoring pages retained to exclude lower-quality material. The resulting dataset sizes are tested in pre-training with the top 40B, 80B, 120B, and 160B tokens. In the initial round, the corpus is limited to the top 40B tokens.

Following the initial round of data gathering, a substantial number of mathematical web pages remain uncollected. This shortfall largely arises because the fastText model was originally trained on a seed corpus of positive samples that lacked adequate heterogeneity. To remedy this, we seek out new mathematical web resources to broaden the variety within the seed dataset, thereby improving the effectiveness of the fastText model. Our procedure starts by segmenting the entire Common Crawl dataset into non-overlapping domains, each defined as a group of web pages sharing a common base URL. We then determine, for every domain, the proportion of web pages included in the first data collection pass. Domains with more than 10\% of their pages collected are labeled as math-oriented (for instance, \url{mathoverflow.net}).

Next, we manually review and label those URLs inside the identified domains that correspond to mathematical content (such as \url{mathoverflow.net/questions}). Web pages under these URLs that have not yet been gathered are appended to the seed corpus. This strategy ensures we obtain more diverse positive samples, which subsequently leads to training a more robust fastText model capable of retrieving additional mathematical web content in future cycles. After executing four data collection cycles, we ultimately compile a set of 35.5 million mathematical web pages, containing a cumulative total of 120 billion tokens. By the fourth cycle, we observe that close to 98\% of the corpus had already been assembled in the prior cycle, prompting us to discontinue further data collection.

To prevent benchmark data from leaking into the training corpus, we adopt the filtering protocol described by \cite{deepseek-coder}: we remove web pages that feature questions or answers present in prominent English mathematics benchmarks such as GSM8K \citep{gsm8k} and MATH \citep{MATH}, as well as Chinese benchmarks like CMATH \citep{wei2023cmath} and AGIEval \citep{agieval}. The filtering methodology is as follows: if any passage of text contains a 10-word string that matches exactly with any substring from these evaluation datasets, we eliminate it from our math corpus. In the case of benchmark texts shorter than 10 grams but at least 3 grams long, we employ an exact match approach to exclude potentially contaminated web content.

\subsection{Validating the Quality of the DeepSeekMath~Corpus}

We perform pre-training experiments to evaluate the DeepSeekMath~Corpus relative to several recently published mathematical training corpora:

\begin{itemize}[topsep=0pt]
    \item \textbf{MathPile} \citep{mathpile}: This dataset aggregates 8.9B tokens of mathematical text from a variety of sources, including textbooks, Wikipedia, ProofWiki, CommonCrawl, StackExchange, and arXiv, with arXiv contributing over 85\% of the content;
    \item \textbf{OpenWebMath} \citep{openwebmath}: Composed of 13.6B tokens, this collection draws from CommonCrawl data filtered specifically for mathematical topics;
    \item \textbf{Proof-Pile-2} \citep{llemma}: This corpus merges OpenWebMath, AlgebraicStack (providing 10.3B tokens of mathematical code), and arXiv preprints (28.0B tokens). For our experiments on Proof-Pile-2, we maintain the arXiv:Web:Code ratio of 2:4:1, as recommended in \cite{llemma}.
\end{itemize}

\subsubsection{Training Setting}
\label{sec:quality-policy}
We conduct mathematical instruction tuning on a pre-existing language model with 1.3 billion parameters, which follows the architectural design of DeepSeek LLMs \citep{deepseek-llm}, referred to here as DeepSeek-LLM 1.3B. Separate models are trained for each mathematical dataset, with each being exposed to 150 billion tokens. The entirety of the training process utilizes the lightweight and resource-efficient HAI-LLM \citep{haillm} framework. In alignment with the DeepSeek LLMs training protocol, we employ the AdamW optimizer \citep{adamW}, using hyperparameters $\beta_1=0.9$, $\beta_2=0.95$, and a weight decay of $0.1$. The learning rate follows a multi-phase schedule: after 2,000 warmup iterations, the rate rises to its maximum value, then drops to 31.6\% of this maximum at the 80\% mark of the total training steps, and subsequently descends further to 10.0\% of the maximum by the 90\% point. The learning rate is capped at 5.3e-4. We use a global batch size of 4 million tokens, with a context window of 4,000 tokens.

\subsubsection{Evaluation Results}

\textbf{The DeepSeekMath~Corpus distinguishes itself by its superior quality, extensive multilingual coverage, and unmatched scale.}

\begin{itemize}[topsep=0pt]
    \item \textbf{High-quality}:
    To assess the effectiveness, we examine downstream outcomes across eight mathematical tasks employing few-shot chain-of-thought prompting \cite{cot}. Results presented in Table \ref{tab:corpora_comparison} indicate that the model fine-tuned on DeepSeekMath~Corpus outperforms others by a substantial margin. Additionally, Figure \ref{fig:corpora_comparisons} reveals that training on DeepSeekMath yields consistently better results than training on Proof-Pile-2 when evaluated at 50B tokens (equivalent to a full pass through Proof-Pile-2), signifying that DeepSeekMath offers higher overall data quality.
    \item \textbf{Multilingual}:
    The DeepSeekMath~Corpus incorporates a variety of languages, with English and Chinese being the most prevalent. According to Table \ref{tab:corpora_comparison}, training on DeepSeekMath~Corpus not only improves mathematical reasoning abilities in English but also considerably enhances Chinese performance. Conversely, currently available math datasets are largely English-focused, resulting in modest gains and even reduced performance on Chinese mathematical reasoning tasks.
    \item \textbf{Large-scale}:
    In terms of volume, DeepSeekMath~Corpus significantly surpasses existing mathematical datasets. Figure \ref{fig:corpora_comparisons} illustrates that DeepSeek-LLM 1.3B benefits from both a sharper improvement trajectory and more persistent performance gains when trained with DeepSeekMath~Corpus. By contrast, the benchmark corpora are much smaller, necessitating repeated epochs within training, which quickly leads to saturated model performance.
\end{itemize}

\subsection{Training and Evaluating DeepSeekMath-Base 7B}

This section details the development of DeepSeekMath-Base 7B, a foundational model characterized by notable reasoning capabilities, with a particular emphasis on mathematical problem solving. The model is initialized from DeepSeek-Coder-Base-v1.5 7B \citep{deepseek-coder}, then trained on a total of 500B tokens. The training data is distributed as follows: DeepSeekMath Corpus accounts for 56\%, AlgebraicStack comprises 4\%, arXiv provides 10\%, 20\% originates from Github code, and the final 10\% consists of natural language data in English and Chinese sourced from Common Crawl. The training methodology largely follows the procedure outlined in Section \ref{sec:quality-policy}, with the exception of setting the maximum learning rate to 4.2e-4 and utilizing a batch size of 10M tokens.

To thoroughly analyze the mathematical proficiency of DeepSeekMath-Base 7B, we examine its competency in generating comprehensive mathematical solutions independently, leveraging external tools for problem-solving, and engaging in formal theorem proving. Additionally, a broader evaluation of the base model is conducted, encompassing natural language understanding, logical reasoning, and programming aptitude.

\paragraph{Mathematical Problem Solving with Step-by-Step Reasoning}

DeepSeekMath-Base’s mathematical problem-solving abilities are assessed using few-shot chain-of-thought prompting \citep{cot} over eight different benchmarks in both English and Chinese. The benchmarks span quantitative reasoning datasets, such as GSM8K \citep{gsm8k}, MATH \citep{MATH}, and CMATH \citep{wei2023cmath}, as well as multiple-choice assessments like MMLU-STEM \citep{mmlu} and Gaokao-MathQA \citep{agieval}. Collectively, these tasks represent a spectrum of mathematical domains, covering levels from elementary education through college curricula.

As indicated in Table \ref{tab:base_math_cot}, DeepSeekMath-Base 7B consistently achieves the highest scores across all eight evaluated benchmarks among open-source base models, encompassing prominent models such as the widely adopted Mistral 7B \citep{mistral} and the newly introduced Llemma 34B \citep{llemma}, both of which have undergone extensive mathematical training with resources like Proof-Pile-2 \citep{llemma}. Particularly striking is DeepSeekMath-Base’s performance on the MATH competition dataset, where it exceeds the best open-source base models by more than 10\% in absolute terms. Furthermore, it even surpasses Minerva 540B \citep{minerva}—a closed-source model that is 77 times larger, built upon the PaLM architecture \citep{palm} and further refined with mathematical texts.

\paragraph{Mathematical Problem Solving with Tool Use} We assess the ability to reason through mathematics with program assistance on GSM8K and MATH using the few-shot program-of-thought prompting approach \citep{pot,pal}. In this setup, models are guided to solve each mathematical problem by producing Python scripts that may use modules like \textit{math} and \textit{sympy} for handling complex calculations. The answer for each task is derived by executing the corresponding program. Table \ref{tab:base_math_pot_proof} demonstrates that DeepSeekMath-Base 7B establishes a new state-of-the-art, outperforming the previous leader, Llemma 34B.

\paragraph{Formal Mathematics}

Automating formal proof construction aids in ensuring the correctness and dependability of mathematical arguments and is attracting growing interest. We measure the performance of DeepSeekMath-Base 7B on informal-to-formal proof generation \citep{dsp_proof}, a task that entails producing a formal proof from a given informal description, along with its formalized statement and accompanying informal proof. The evaluation employs the miniF2F benchmark \citep{minif2f}, targeting formal Olympiad-level problems; for each, the model generates an Isabelle formal proof using few-shot prompting. Consistent with \cite{dsp_proof}, we have the model output an outline of the proof and subsequently apply the Sledgehammer automated theorem prover \citep{sledgehammer} to automatically fill in omitted steps. Results, as presented in Table \ref{tab:base_math_pot_proof}, confirm the strong capabilities of DeepSeekMath-Base 7B in automating the formalization of mathematical proofs.

\paragraph{Natural Language Understanding, Reasoning, and Code}

We further evaluate the model's proficiency in natural language understanding with MMLU \citep{mmlu}, logical reasoning on BBH \citep{bbh}, and programming competence on both HumanEval \citep{codex} and MBPP \citep{mbpp}. Table \ref{tab:understanding_reasoning_code} indicates that DeepSeekMath-Base 7B delivers marked improvements over its predecessor, DeepSeek-Coder-Base-v1.5 \citep{deepseek-coder}, on MMLU and BBH, highlighting the beneficial effect of mathematical training on language and reasoning abilities. Additionally, by incorporating coding tokens during continued training, DeepSeekMath-Base 7B retains the strong coding performance exhibited by DeepSeek-Coder-Base-v1.5 on HumanEval and MBPP. Overall, DeepSeekMath-Base 7B significantly outperforms the general-purpose Mistral 7B \citep{mistral} on all three benchmarks concerning reasoning and coding.

\section{Supervised Fine-Tuning}

\subsection{SFT Data Curation}

We compile an instruction-tuning dataset for mathematics, comprising both English and Chinese problems sourced from a variety of mathematical domains and difficulty levels. Each problem is paired with a corresponding solution formatted as chain-of-thought (CoT) \citep{cot}, program-of-thought (PoT) \citep{pot,pal}, or tool-integrated reasoning \citep{tora}. The dataset contains a total of 776,000 training instances.

\begin{itemize}[topsep=0pt]
    \item \textbf{English mathematical datasets}: We enhance GSM8K and MATH questions with tool-integrated answer annotations. In addition, we utilize selected samples from MathInstruct \citep{MathInstruct}, as well as the training portion of Lila-OOD \citep{lila}, in which questions are addressed through CoT or PoT strategies. The English dataset encompasses a broad spectrum of mathematics, including areas such as algebra, probability, number theory, calculus, and geometry.
    \item \textbf{Chinese mathematical datasets}: Our collection consists of Chinese K-12 mathematics problems covering 76 detailed sub-topics, such as linear equations. Solutions for these problems are annotated in both CoT and tool-integrated reasoning formats.
\end{itemize}

\subsection{Training and Evaluating DeepSeekMath-Instruct 7B}

Here, we present the procedure for mathematical instruction tuning of DeepSeekMath-Instruct 7B, which is fine-tuned from DeepSeekMath-Base. For model training, examples are concatenated randomly up to a maximum of 4,000 tokens per input. The model is trained for 500 iterations using a batch size of 256 and a fixed learning rate set at 5e-5.

The model's ability to handle mathematical tasks—both with and without external tool usage—is assessed on four benchmarks targeting quantitative reasoning in English and Chinese. We compare our approach with contemporary state-of-the-art systems:

\begin{itemize}[topsep=0pt]
    \item \textbf{Closed-source models} examined include: (1) the GPT series, with GPT-4 \citep{gpt4} and GPT-4 Code Interpreter~\footnote{\url{https://openai.com/blog/chatgpt-plugins\#\#code-interpreter}} considered most capable; (2) Gemini Ultra and Pro \citep{gemini}; (3) Inflection-2 \citep{inflection-2}; (4) Grok-1~\footnote{\url{https://x.ai/model-card}}; and (5) models recently introduced by Chinese technology firms, such as Baichuan-3~\footnote{\url{https://www.baichuan-ai.com}}, and (6) the new GLM-4~\footnote{\url{https://open.bigmodel.cn/dev/api\#glm-4}} from the GLM series \citep{glm}.
\end{itemize}

These models are designed for broad applicability, with the majority having completed various alignment processes. 
\item \textbf{Open-source models} span: general-purpose models such as (1) DeepSeek-LLM-Chat 67B \citep{deepseek-llm}, (2) Qwen 72B \citep{qwen}, (3) SeaLLM-v2 7B \citep{seallm}, and (4) ChatGLM3 6B \citep{chatglm3}; as well as models specialized for mathematical capabilities, including (5) InternLM2-Math 20B\footnote{\url{https://github.com/InternLM/InternLM-Math}}, which extends InternLM2 by incorporating mathematical pretraining followed by instruction tuning, (6) Math-Shepherd-Mistral 7B, which applies PPO training \citep{schulman2017proximal} on Mistral 7B \citep{mistral} with a process-supervised reward framework, (7) the WizardMath series \citep{wizardmath}, aimed at enhancing mathematical reasoning in Mistral 7B and Llama-2 70B \citep{llama2} using an evolve-instruct method (a form of instruction tuning employing AI-generated instructions) and PPO on datasets primarily consisting of GSM8K and MATH problems, (8) MetaMath 70B \citep{metamath}, which fine-tunes Llama-2 70B on an expanded set of GSM8K and MATH, (9) ToRA 34B \cite{tora}, a version of CodeLlama 34B trained to carry out tool-integrated mathematical reasoning, and (10) MAmmoTH 70B \citep{MathInstruct}, which involves Llama-2 70B tuned on MathInstruct through instruction tuning.
\end{itemize}

According to Table \ref{tab:sft_rl_math}, in the context where the use of external tools is not permitted, DeepSeekMath-Instruct 7B exhibits robust step-by-step reasoning performance. Particularly on the challenging MATH dataset, this model outperforms all publicly available open-source models, as well as the majority of closed-source models (including Inflection-2 and Gemini Pro), with a margin of at least 9\% in absolute terms. This result holds even in comparison to models that possess much larger parameter counts (for instance, Qwen 72B) or have benefited from math-specific reinforcement learning enhancements (such as WizardMath-v1.1 7B). While DeepSeekMath-Instruct matches the performance of leading Chinese proprietary models GLM-4 and Baichuan-3 on the MATH benchmark, it is still exceeded by the likes of GPT-4 and Gemini Ultra.

When evaluated under conditions that allow models to leverage both natural language reasoning and programmatic tools for solving problems, DeepSeekMath-Instruct 7B reaches nearly 60\% accuracy on the MATH dataset, outperforming all prior open-source competitors. On other standard benchmarks, it achieves results on par with DeepSeek-LLM-Chat 67B, which is an order of magnitude larger in scale.

\section{Reinforcement Learning}

\subsection{Group Relative Policy Optimization} 
Reinforcement learning (RL) has been shown to effectively further strengthen the mathematical reasoning capacities of large language models after completing Supervised Fine-Tuning (SFT) \citep{wang2023math, wizardmath}. In the subsequent discussion, we present Group Relative Policy Optimization (GRPO), our efficient and effective RL technique.

\subsubsection{From PPO to GRPO} 
Proximal Policy Optimization (PPO) \citep{schulman2017proximal} stands as a prevalent actor-critic reinforcement learning approach utilized during the RL fine-tuning stage of large language models \citep{ouyang2022training}. The approach focuses on maximizing the following surrogate loss function:

\begin{equation}
\footnotesize
    \mathcal{J}_{PPO}(\theta) = \mathbb{E}{[q \sim P(Q), o \sim \pi_{\theta_{old}}(O|q)]} \frac{1}{|o|} \sum_{t=1}^{|o|} \min \left[ \frac{\pi_\theta(o_{t} | q, o_{<t})}{\pi_{\theta_{old}}(o_{t} | q, o_{<t})} A_{t}, \text{clip} \left( \frac{\pi_\theta(o_{t} | q, o_{<t})}{\pi_{\theta_{old}}(o_{t} | q,

Here, $\pi_{\theta}$ denotes the current policy model, while $\pi_{\theta_{old}}$ represents the previous policy. The variables $q$ and $o$ refer to questions and outputs, which are drawn from the question dataset and the former policy $\pi_{\theta_{old}}$, respectively. The parameter $\varepsilon$ is a hyper-parameter associated with clipping in PPO, intended to stabilize the training procedure. The term $A_t$ stands for the advantage, calculated through Generalized Advantage Estimation (GAE) \citep{gae}, based on subsequent rewards $\{r_{\ge t}\}$ as well as a learned value function $V_{\psi}$. Therefore, within PPO, it is necessary to jointly optimize a value function alongside the main policy model. To avoid excessive reward model optimization, a widely-adopted strategy is to include a per-token KL divergence penalty between a reference model and the policy in the reward at each token \citep{ouyang2022training}, as specified by

\begin{equation}
    r_{t} = r_\varphi(q, o_{\le t}) - \beta \log\frac{\pi_{\theta}(o_{t}|q, o_{<t})}{\pi_{ref}(o_{t}|q, o_{<t})},
\label{eq:PPO-reward}
\end{equation}

In this expression, $r_\varphi$ denotes the reward model, $\pi_{ref}$ refers to the reference model (often initialized as the SFT model), and $\beta$ regulates the strength of the KL penalty.

Because the value function in PPO is generally parameterized as another large neural model similar in size to the policy model, this dual-model setup leads to significant memory consumption and computational demands. In addition, during reinforcement learning, the value function acts as a baseline in advantage estimation, contributing to variance reduction. However, in large language model scenarios, the reward model is typically applied only to the last token, making it challenging to obtain precise value function targets across all token positions. To mitigate these challenges, as illustrated in Figure \ref{fig:grpo}, we introduce Group Relative Policy Optimization (GRPO), a method that eliminates the need for an explicit value function estimator like that used in PPO. Instead, GRPO adopts the mean reward across a collection of sampled outputs—each responding to the same input question—as a baseline. Specifically, for a given question $q$, GRPO samples a set of outputs $\{o_1, o_2, \cdots, o_G\}$ from the previous policy $\pi_{\theta_{old}}$, and then maximizes the policy objective:

\begin{equation}
\footnotesize
\begin{split}
    \mathcal{J}_{GRPO}(\theta) &= \mathbb{E}{[q \sim P(Q), \{o_i\}_{i=1}^G \sim \pi_{\theta_{old}}(O|q)]}  \\
    & \frac{1}{G}\sum_{i=1}^G\frac{1}{|o_i|} \sum_{t=1}^{|o_i|} \left\{ \min \left[ \frac{\pi_\theta(o_{i,t} | q, o_{i,<t})}{\pi_{\theta_{old}}(o_{i,t} | q, o_{i,<t})} \hat{A}_{i,t}, \text{clip} \left( \frac{\pi_\theta(o_{i,t} | q, o_{i,<t})}{\pi_{\theta_{old}}(o_{i,t} | q, o_{i,<t})}, 1 - \varepsilon, 1 + \varepsilon \right)  \hat{A}_{i,t} \right] - \beta \mathbb{D}_{KL}\left[\pi_{\theta} || \pi_{ref}\right]\right\} ,
\end{split}
\label{eq:GRPO-obj}
\end{equation}

with $\varepsilon$ and $\beta$ acting as hyper-parameters, and where $\hat{A}_{i,t}$ refers to an advantage term computed solely from the relative rewards of group outputs—details of which will be elaborated in subsequent sections. The group-based approach for calculating advantages in GRPO is particularly suited to the typical design of reward models, which are trained on comparative annotations among different outputs for the same input question. Furthermore, GRPO regularizes the policy by directly adding the KL divergence between the updated policy and the reference policy as a term in the objective function, in contrast to the approach of incorporating KL penalty within the reward as in PPO. This avoids adding complexity to the calculation of $\hat{A

\begin{algorithm}[t]
  \small
  \caption{Iterative Group Relative Policy Optimization}
  \textbf{Input} initial policy model $\pi_{\theta_{\text{init}}}$; reward models $r_\varphi$; task prompts $\mathcal{D}$; 
  hyperparameters $\varepsilon$, $\beta$, $\mu$
  \begin{algorithmic}[1]
    \State Initialize policy model $\pi_\theta \leftarrow \pi_{\theta_{\text{init}}}$
    \For{iteration = 1, \dots, I}
       \State Set reference model $\pi_{ref} \leftarrow \pi_{\theta}$
      \For{step = 1, \dots, M}
      \State Draw a mini-batch $\mathcal{D}_b$ from $\mathcal{D}$
      \State Assign current policy to $\pi_{\theta_{old}} \leftarrow \pi_{\theta}$ 
      \State For each question $q \in \mathcal{D}_b$, generate $G$ outputs $\{o_i\}_{i=1}^G \sim \pi_{\theta_{old}} (\cdot \mid q) $
      \State Evaluate sampled outputs $\{o_i\}$ using $r_\varphi$ to obtain corresponding rewards $\{r_i\}_{i=1}^{G}$
      \State Employ group relative advantage estimation to calculate $\hat{A}_{i,t}$ for the $t$-th token of each $o_i$
      \For{GRPO iteration = 1, \dots, $\mu$}
        \State Update policy $\pi_{\theta}$ via maximization of the GRPO objective (Equation \ref{eq:GC-GRPO})
      \EndFor
    \EndFor 
    \State Continuously retrain $r_\varphi$ by leveraging a replay buffer mechanism.
    \EndFor 
  \end{algorithmic}
  \textbf{Output} $\pi_\theta$
  \label{alg:iter-grpo}
\end{algorithm}

\subsubsection{Outcome Supervision RL with GRPO} Consider a scenario in which, for every question $q$, a set of $G$ outputs $\{o_1, o_2, \cdots, o_G\}$ are produced using the previous policy model $\pi_{\theta_{old}}$. Each of these outputs is then assessed by a reward model, resulting in a set of reward values $\mathbf{r}=\{r_1, r_2, \cdots, r_G\}$. The obtained rewards are standardized: the group mean is subtracted from each reward and then divided by the group's standard deviation. Under outcome supervision, this normalized reward is provided only at the end of every output $o_i$ and is uniformly assigned as the advantage $\hat{A}_{i, t}$ to every token in $o_i$, i.e., $\hat{A}_{i, t} = \widetilde{r}_i = \frac{r_i- {\rm mean}(\mathbf{r})}{{\rm std}(\mathbf{r})}$. Optimization then proceeds via maximization of the objective detailed in equation (\ref{eq:GRPO-obj}).

\subsubsection{Process Supervision RL with GRPO} While outcome supervision offers terminal rewards, this approach may not suffice for guiding the model in intricate mathematical problem solving. Following \cite{wang2023math}, process supervision is incorporated, delivering rewards after each intermediate reasoning step. Specifically, given a prompt $q$ and $G$ sampled sequences $\{o_1, o_2, \cdots, o_G\}$, a process-level reward model is used to provide a reward for every step, denoted as: $\mathbf{R} = \{ \{r_1^{index(1)},\cdots,r_1^{index(K_1)}\}, \cdots,  \{r_G^{index(1)},\cdots,r_G^{index(K_G)}\} \}$, where $index(j)$ refers to the position of the $j$-th step's final token and $K_i$ indicates how many steps are in output $i$. Each of these stepwise rewards is standardized via $\widetilde{r}_i^{index(j)} = \frac{r_i^{index(j)} - {\rm mean(\mathbf{R})}}{{\rm std(\mathbf{R})}}$. Afterwards, the advantage assigned to every token is computed by summing the normalized rewards from subsequent steps, formally, $\hat{A}_{i, t} = \sum_{index(j) \ge t} \widetilde{r}_i^{index(j)}$.

\subsubsection{Iterative RL with GRPO} As the RL training advances, previous versions of the reward model may no longer provide adequate guidance for the evolving policy model. To address this, we investigate iterative RL with GRPO. As detailed in Algorithm \ref{alg:iter-grpo}, this approach involves using the outputs from the policy model to generate updated datasets for reward model training. The reward model is then continually trained using a replay strategy that maintains 10\% of prior data. Subsequently, the policy model serves as the new reference model, and policy updates are carried out using the improved reward model.

\subsection{Training and Evaluating DeepSeekMath-RL}

Our RL procedures utilize DeepSeekMath-Instruct 7B as the base policy model. For RL training, we restrict the dataset to chain-of-thought-style questions derived from GSM8K and MATH within the SFT data, totaling about 144K questions. Questions from other SFT sources are omitted in order to specifically measure RL’s effects on benchmarks without additional RL-stage data exposure. The reward model’s training set is curated according to the methodology from \citep{wang2023math}. The initial reward model is trained on DeepSeekMath-Base 7B with a learning rate of 2e-5. In the GRPO procedure, the policy model’s learning rate is set to 1e-6 and the KL divergence coefficient is 0.04. For each prompt, $64$ candidate responses are generated. Maximum sequence length and batch size are both set at 1024. The policy model undergoes one update per exploration iteration. Evaluation of DeepSeekMath-RL 7B is performed using the same benchmarks as for DeepSeekMath-Instruct 7B. Within this setup, tasks from GSM8K and MATH using chain-of-thought reasoning are considered in-domain, while all other evaluation benchmarks are treated as out-of-domain.

Table \ref{tab:sft_rl_math} presents a comparative analysis of open- and closed-source models, evaluating both their chain-of-thought and tool-integrated reasoning abilities on English and Chinese datasets. Our observations are as follows: 1) When utilizing chain-of-thought reasoning, DeepSeekMath-RL 7B achieves accuracy rates of 88.2\% on GSM8K and 51.7\% on MATH, outperforming every open-source model in the 7B–70B category and exceeding most closed-source alternatives as well. 2) Notably, the only additional supervision for DeepSeekMath-RL 7B comes from instruction tuning data in the chain-of-thought format for GSM8K and MATH, initialized from DeepSeekMath-Instruct 7B. Even though its supervised data is limited in scope, DeepSeekMath-RL 7B consistently surpasses DeepSeekMath-Instruct 7B on every evaluation metric, highlighting the efficacy of reinforcement learning.

\section{Discussion}
Here, we discuss our empirical results from both pre-training and reinforcement learning experiments.

\subsection{Insights from Pre-Training}

We begin with our key insights regarding pre-training. Unless noted otherwise, all results conform to the training procedures specified in Section \ref{sec:quality-policy}. The term DeepSeekMath Corpus, as used throughout this section, refers to the 89B-token dataset compiled during the second data collection phase.

\subsubsection{Code Training Enhances Mathematical Reasoning}

There is a widely held—though previously untested—assumption that training on code data enhances a model’s reasoning capabilities. Our investigation provides partial support for this claim, specifically within mathematics: integrating code data into training bolsters the model's mathematical reasoning, both with and without external tool assistance.

To examine the impact of code training on mathematical reasoning, we carried out experiments contrasting a two-stage training regimen with a one-stage training paradigm:

\textbf{Two-Stage Training}

\begin{itemize}[topsep=0pt]
    \item \textbf{Code Training for 400B Tokens $\rightarrow$ Math Training for 150B Tokens}: DeepSeek-LLM 1.3B is first exposed to 400B tokens of code data, subsequently followed by an additional 150B tokens of mathematics data for continued training;
    \item \textbf{General Training for 400B Tokens $\rightarrow$ Math Training for 150B Tokens}: For comparison, we perform an alternative experiment using 400B general tokens (drawn from DeepSeek-AI’s extensive general corpus) instead of code in the initial training stage, thereby enabling analysis of whether code pre-training offers specific benefits for mathematical reasoning over general domain pre-training.
\end{itemize}

\textbf{One-Stage Training}

\begin{itemize}[topsep=0pt]
    \item \textbf{Math Training for 150B Tokens}: DeepSeek-LLM 1.3B undergoes training exclusively on 150B mathematics tokens;
    \item \textbf{Training on a mixture of 400B Code Tokens and 150B Math Tokens}: Sequential math training after code exposure diminishes code abilities. Thus, we assess whether jointly presenting code and math tokens during a single training phase can both enhance mathematical reasoning and help preserve prior coding skills by alleviating catastrophic forgetting.
\end{itemize}

\paragraph{Results}
The performance results across the varied training configurations are summarized in Table \ref{tab:code-math} and Table \ref{tab:code-math-general-eval}.

Introducing code data into training leads to improved program-aided mathematical reasoning, as observed for both staged and simultaneous training paradigms. Specifically, in the two-stage setting presented in Table \ref{tab:code-math}, pretraining with code tokens alone produces notable gains in solving tasks such as GSM8K and MATH using Python. Applying math-specific training thereafter further enhances results. Notably, simultaneous exposure to mixed code and math data in one-stage training reduces the extent of catastrophic forgetting characteristic of staged training, while also jointly bolstering coding (Table \ref{tab:code-math-general-eval}) and program-supported mathematical reasoning (Table \ref{tab:code-math}).

Furthermore, code pretraining contributes positively to mathematical reasoning capabilities even when no tool use is involved. In the two-stage scheme, the code phase alone delivers a moderate boost and serves to accelerate progress during the following math-focused phase, ultimately producing optimal results. Conversely, mixing code and math tokens in a single training phase negatively affects mathematical reasoning without tool support. A possible explanation is that the relatively small scale of DeepSeek-LLM 1.3B restricts its ability to simultaneously master the intricacies of both programming and mathematical content during joint training.

\subsubsection{ArXiv Papers Show Little Promise for Enhancing Mathematical Reasoning}

While arXiv papers are routinely included as part of math pre-training datasets \citep{minerva,gpt-f,llemma,mathpile}, their specific influence on mathematical reasoning remains largely unexplored. Somewhat unexpectedly, our findings indicate that pre-training on arXiv material does not meaningfully bolster mathematical reasoning abilities. We evaluate this by training models of distinct scales—specifically, DeepSeek-LLM 1.3B and DeepSeek-Coder-Base-v1.5 7B \citep{deepseek-coder}—on different versions of arXiv datasets subjected to unique preprocessing schemes:

\begin{itemize}[topsep=0pt]
    \item \textbf{MathPile} \citep{mathpile}: An 8.9-billion-token collection curated with filtering and heuristic cleaning, composed of over 85\% scientific articles from arXiv;
    \item \textbf{ArXiv-RedPajama} \citep{redpajama}: A 28-billion-token dump of arXiv LaTeX documents with extraneous material such as preambles, comments, macros, and references removed.
\end{itemize}

In our trials, DeepSeek-LLM 1.3B was trained for 150B tokens and DeepSeek-Coder-Base-v1.5 7B for 40B tokens on each version of the arXiv data. Models pre-trained solely on these arXiv corpora demonstrate no clear gains and sometimes even declines in performance across an array of mathematical evaluation sets of varying difficulty. These evaluation sets comprise quantitative reasoning benchmarks like GSM8K and MATH (see Table \ref{tab:arxiv-cot}), STEM-focused multiple-choice tasks as in MMLU-STEM (refer to Table \ref{tab:arxiv-cot}), as well as formal mathematics challenges like miniF2F (see Table \ref{tab:arxiv_atp}).

Nevertheless, the scope of this conclusion is limited and should be considered cautiously. In particular, our analysis does not yet address:

\begin{itemize}[topsep=0pt]
    \item The effect of arXiv pre-training on other math-related capabilities absent from this work, such as translating formal mathematical statements or proofs into informal language;
    \item The interaction between arXiv tokens and other datasets during joint training;
    \item Whether larger-scale models could derive more substantial benefit from arXiv-based corpora.
\end{itemize}

Accordingly, these open questions highlight the need for additional investigation, which we defer to future work.

\subsection{Reinforcement Learning Insights}
\subsubsection{Toward a Unified Training Paradigm} In this segment, we develop a cohesive framework to analyze a variety of training procedures, including SFT, RFT, DPO, PPO, GRPO, and further investigate how the framework’s elements influence outcomes. Conventionally, the parameter gradient for a training objective can be formalized as follows:

\begin{equation}
    \nabla_{\theta}\mathcal{J}_{\textcolor{red}{\mathcal{A}}}(\theta) = \mathbb{E}[\underbrace{(q,o) \sim \textcolor{red}{\mathcal{D}}}_{Data \ Source}]\left( \frac{1}{|o|} \sum_{t=1}^{|o|}  \underbrace{GC_{{\mathcal{A}}}(q, o, t, \textcolor{red}{\pi_{{rf}}})}_{Gradient \ Coefficient}  \nabla_{\theta}\log \pi_{\theta}(o_t | q, o_{<t})\right).
\label{eq:objective}
\end{equation}

This framework comprises three principal elements: 1) the \textit{Data Source $\mathcal{D}$}, which supplies the examples used during training; 2) the \textit{Reward Function $\pi_{{rf}}$}, providing the evaluative feedback signal for learning; and 3) the \textit{Algorithm $\mathcal{A}$}, which integrates both the data and reward information to compute the gradient coefficient $GC$ that modulates the degree of updating—be it punitive or incentivizing—applied to the training samples. Within this unified formulation, several commonly used approaches can be delineated:

\begin{itemize}
    \item  \textbf{Supervised Fine-tuning (SFT)}: This approach adapts a pretrained model using a dataset curated by humans specifically for SFT purposes.
    \item \textbf{Rejection Sampling Fine-tuning (RFT)}: Here, the SFT model undergoes additional fine-tuning with outputs drawn from the SFT model and further screened for correctness according to SFT queries.
    \item \textbf{Direct Preference Optimization (DPO)}: DPO takes the SFT model as a starting point and continues fine-tuning by leveraging additional outputs sampled from the SFT model, employing a pairwise loss function specific to DPO.
    \item \textbf{Online Rejection Sampling Fine-tuning (Online RFT)}: Distinct from RFT, Online RFT starts with the SFT model as the initial policy but subsequently updates using augmented outputs that are generated from the policy model in its current state, thus allowing for continual adaptation.
    \item \textbf{PPO/GRPO}: These methods initialize the policy with the SFT model and apply reinforcement via outputs sampled directly from the live policy during training.
\end{itemize}

A summary of the elements comprising these approaches is presented in Table \ref{tab:data_policy}. For an expanded explanation of the derivation steps, consult Appendix \ref{app:analysis-rl}.

\paragraph{Discussion on Data Source} We classify data sources into two main types: online and offline sampling. Online sampling refers to generating training data by exploring outcomes of the current, actively learning policy model, whereas offline sampling uses data produced from the initial SFT model’s outputs. RFT and DPO both operate in the offline setting, whereas Online RFT and GRPO adopt online sampling procedures.

As depicted in Figure \ref{fig:rl-analysis}, our results indicate that Online RFT demonstrates notable improvements over RFT across two benchmark tasks. Although performance between Online RFT and RFT remains similar in the early phases of training, Online RFT establishes a distinct lead as training advances, underscoring the benefits of online data collection. This observation aligns with expectations: initially, there is considerable similarity between the actor and the SFT model, so their sampled outputs differ only marginally. At later stages, however, actor-generated samples diverge more substantially, and dynamically collected training data proves increasingly advantageous.

\paragraph{Discussion on Gradient Coefficient} To update model parameters, each algorithm translates the reward signal into a gradient coefficient. In our analysis, the reward functions are categorized as either `Rule' or `Model'. `Rule' entails evaluating responses strictly according to answer correctness, whereas `Model' involves constructing a reward model—trained on rule-based assessments—to assign scores to responses. Notably, Equations \ref{eq:GC-RFT} and \ref{eq:GC-GRPO} reveal a critical distinction between GRPO and Online RFT: GRPO uniquely modulates the gradient coefficient according to the reward values estimated by the reward model, thereby facilitating a range of positive and negative reinforcements dependent on the response’s quality. By contrast, Online RFT treats all responses with correct answers uniformly, lacking explicit penalization for incorrect ones and applying consistent reinforcement to correct responses.

As illustrated in Figure \ref{fig:rl-analysis}, GRPO outperforms online RFT, underscoring the efficacy of modifying the positive and negative gradient coefficients. Moreover, the results reveal that GRPO+PS outshines GRPO+OS, suggesting that utilizing more granular, step-sensitive gradient coefficients offers notable advantages. Additionally, we investigate the impact of iterative RL, conducting two consecutive training rounds in our experiments. The results presented in Figure \ref{fig:iter-rl} indicate that iterative RL notably boosts performance, with the most pronounced gains observed after the first round.

\subsubsection{Why RL Works?} 
In this study, reinforcement learning is applied to a subset of the instruction tuning dataset, yielding substantial improvements over the instruction-tuned baseline model. To provide deeper insight into the mechanisms underlying these gains, we measure both Pass@K and Maj@K accuracies for the Instruct and RL-optimized models across two evaluation benchmarks. As depicted in Figure \ref{fig:maj-pass}, while RL substantially elevates Maj@K accuracy, it does not significantly impact Pass@K. These results imply that RL’s benefits stem from a more robust output distribution—specifically, \textbf{the primary advancement lies in amplifying the likelihood of correct answers among the TopK outputs, rather than augmenting the model’s underlying abilities}. This pattern is consistent with findings from \citep{wang2023making}, which describe a \textbf{misalignment problem} in reasoning tasks within SFT models and demonstrate that SFT reasoning performance can be bolstered via various preference alignment approaches \citep{yuan2023rrhf,song2023preference,wang2023making}.

\subsubsection{How to Achieve More Effective RL?}
\label{sec:effec_rl}
Our experiments demonstrate that reinforcement learning is particularly effective for mathematical reasoning tasks. Furthermore, we propose a unified perspective that encompasses various major training frameworks, interpreting all of them as forms or simplifications of RL methodologies. As formalized in Equation \ref{eq:objective}, these approaches consist of three fundamental elements: Data Source, Algorithm, and Reward Function. We discuss promising avenues for development within each of these aspects.

\paragraph{Data Source} The data source serves as the foundation for all training strategies. In the RL framework, it specifically refers to unlabeled problems and the corresponding outputs sampled from the policy model. In this work, our experiments utilize only questions originating from the instruction tuning process, paired with basic nucleus sampling for generating outputs. We speculate this choice partly explains why our RL process yields improvements mainly in Maj@K. For future work, we plan to assess the RL method on prompts beyond the distribution of the training set and explore the use of \textbf{advanced sampling (decoding) methods}, such as those leveraging tree-search techniques \citep{yao2023tree}. In addition, integrating \textbf{efficient inference strategies} \citep{xia-etal-2023-speculative,leviathan2023fast,kwon2023efficient,xia2024unlocking}, which critically affect how effectively the policy model can explore its output space, will also be an essential consideration.

\paragraph{Algorithms} Algorithms utilize the data alongside reward signals to generate gradient coefficients, which are then employed to adjust the model parameters. According to Equation \ref{eq:objective}, most contemporary techniques fundamentally \textbf{TRUST} the reward function to modulate the conditional likelihood of a given token, either increasing or decreasing it accordingly. Nevertheless, it is not feasible to guarantee the consistent reliability of the reward signal, particularly when dealing with tasks of high complexity. For instance, even within the meticulously curated PRM800K datasets \citep{lightman2023let}, assembled by skilled annotators, there remains a notable portion—approximately 20\%—of erroneous annotations\footnote{\url{https://github.com/openai/prm800k/issues/12\#issuecomment-1728491852}}. Therefore, it is essential to investigate reinforcement learning strategies that maintain robustness in the presence of noisy reward feedback. We propose that alignment frameworks inspired by \textbf{WEAK-TO-STRONG} \citep{burns2023weak} could fundamentally reshape existing learning methodologies.

\paragraph{Reward Function} The reward function supplies the principal signal that guides training. In reinforcement learning, this reward is typically defined by a neural reward model. We identify three key avenues for advancing reward models: 1) \textbf{Improving the generalization capacity of reward models.} Reward models must extend their predictive power to address out-of-distribution queries and complex model generations; failing this, reinforcement learning may serve only to reinforce the language model's existing behavior without driving significant improvement; 2) \textbf{Capturing the uncertainty associated with reward models.} Quantifying uncertainty can serve as a critical interface between suboptimal reward signals and robust weak-to-strong alignment algorithms; 3) \textbf{Developing high-quality process-based reward models efficiently} so as to deliver more nuanced, granular training signals during multi-step reasoning tasks \citep{lightman2023let,wang2023math}.

\section{Conclusion, Limitation, and Future Work}

This paper introduces DeepSeekMath, an open-source model that establishes state-of-the-art results on the competition-grade MATH benchmark, narrowing the gap with proprietary models. DeepSeekMath~originates from the DeepSeek-Coder-v1.5 7B model, followed by continual training over 500B tokens, where 120B math-specific tokens are obtained from Common Crawl and comprise a significant proportion of the training data. Our comprehensive ablation analysis indicates that web sources provide abundant, high-quality mathematical content, while the benefit from arXiv sources falls short of our expectations. We present Group Relative Policy Optimization (GRPO), an adaptation of Proximal Policy Optimization (PPO), which substantially enhances mathematical reasoning abilities while requiring less memory. Experimental evidence demonstrates the efficacy of GRPO, even when DeepSeekMath-Instruct 7B attains strong benchmark results. We further offer a unified framework for interpreting a range of reinforcement learning techniques and propose several promising avenues for advancing reinforcement learning effectiveness.

While DeepSeekMath~exhibits outstanding performance on quantitative reasoning challenges, it still lags behind proprietary systems in geometry and theorem-proving tasks. For example, preliminary evaluations reveal that the model struggles with geometric problems involving triangles and ellipses, suggesting possible biases in both pre-training and fine-tuning data selection. Additionally, due to its scale, DeepSeekMath~does not match GPT-4 in few-shot settings; whereas GPT-4 demonstrates clear improvements from few-shot prompts, DeepSeekMath~performs comparably under zero-shot and few-shot scenarios. Going forward, our plans include refining our data selection strategies to create even higher quality pre-training corpora. Moreover, we aim to investigate the research paths highlighted in Section \ref{sec:effec_rl} to further enhance the reinforcement learning capabilities of large language models (LLMs).